\IfFileExists{stacks-project.cls}{%
\documentclass{stacks-project}
}{%
\documentclass{amsart}
}

% For dealing with references we use the comment environment
\usepackage{verbatim}
\newenvironment{reference}{\comment}{\endcomment}
%\newenvironment{reference}{}{}
\newenvironment{slogan}{\comment}{\endcomment}
\newenvironment{history}{\comment}{\endcomment}

% For commutative diagrams we use Xy-pic
\usepackage[all]{xy}

% We use 2cell for 2-commutative diagrams.
\xyoption{2cell}
\UseAllTwocells

% We use multicol for the list of chapters between chapters
\usepackage{multicol}

% This is generally recommended for better output
\usepackage{lmodern}
\usepackage[T1]{fontenc}

% For cross-file-references
\usepackage{xr-hyper}

% Package for hypertext links:
\usepackage{hyperref}

% For any local file, say "hello.tex" you want to link to please
% use \externaldocument[hello-]{hello}
\externaldocument[introduction-]{introduction}
\externaldocument[conventions-]{conventions}
\externaldocument[sets-]{sets}
\externaldocument[categories-]{categories}
\externaldocument[topology-]{topology}
\externaldocument[sheaves-]{sheaves}
\externaldocument[sites-]{sites}
\externaldocument[stacks-]{stacks}
\externaldocument[fields-]{fields}
\externaldocument[algebra-]{algebra}
\externaldocument[brauer-]{brauer}
\externaldocument[homology-]{homology}
\externaldocument[derived-]{derived}
\externaldocument[simplicial-]{simplicial}
\externaldocument[more-algebra-]{more-algebra}
\externaldocument[smoothing-]{smoothing}
\externaldocument[modules-]{modules}
\externaldocument[sites-modules-]{sites-modules}
\externaldocument[injectives-]{injectives}
\externaldocument[cohomology-]{cohomology}
\externaldocument[sites-cohomology-]{sites-cohomology}
\externaldocument[dga-]{dga}
\externaldocument[dpa-]{dpa}
\externaldocument[sdga-]{sdga}
\externaldocument[hypercovering-]{hypercovering}
\externaldocument[schemes-]{schemes}
\externaldocument[constructions-]{constructions}
\externaldocument[properties-]{properties}
\externaldocument[morphisms-]{morphisms}
\externaldocument[coherent-]{coherent}
\externaldocument[divisors-]{divisors}
\externaldocument[limits-]{limits}
\externaldocument[varieties-]{varieties}
\externaldocument[topologies-]{topologies}
\externaldocument[descent-]{descent}
\externaldocument[perfect-]{perfect}
\externaldocument[more-morphisms-]{more-morphisms}
\externaldocument[flat-]{flat}
\externaldocument[groupoids-]{groupoids}
\externaldocument[more-groupoids-]{more-groupoids}
\externaldocument[etale-]{etale}
\externaldocument[chow-]{chow}
\externaldocument[intersection-]{intersection}
\externaldocument[pic-]{pic}
\externaldocument[weil-]{weil}
\externaldocument[adequate-]{adequate}
\externaldocument[dualizing-]{dualizing}
\externaldocument[duality-]{duality}
\externaldocument[discriminant-]{discriminant}
\externaldocument[derham-]{derham}
\externaldocument[local-cohomology-]{local-cohomology}
\externaldocument[algebraization-]{algebraization}
\externaldocument[curves-]{curves}
\externaldocument[resolve-]{resolve}
\externaldocument[models-]{models}
\externaldocument[functors-]{functors}
\externaldocument[equiv-]{equiv}
\externaldocument[pione-]{pione}
\externaldocument[etale-cohomology-]{etale-cohomology}
\externaldocument[proetale-]{proetale}
\externaldocument[relative-cycles-]{relative-cycles}
\externaldocument[more-etale-]{more-etale}
\externaldocument[trace-]{trace}
\externaldocument[crystalline-]{crystalline}
\externaldocument[spaces-]{spaces}
\externaldocument[spaces-properties-]{spaces-properties}
\externaldocument[spaces-morphisms-]{spaces-morphisms}
\externaldocument[decent-spaces-]{decent-spaces}
\externaldocument[spaces-cohomology-]{spaces-cohomology}
\externaldocument[spaces-limits-]{spaces-limits}
\externaldocument[spaces-divisors-]{spaces-divisors}
\externaldocument[spaces-over-fields-]{spaces-over-fields}
\externaldocument[spaces-topologies-]{spaces-topologies}
\externaldocument[spaces-descent-]{spaces-descent}
\externaldocument[spaces-perfect-]{spaces-perfect}
\externaldocument[spaces-more-morphisms-]{spaces-more-morphisms}
\externaldocument[spaces-flat-]{spaces-flat}
\externaldocument[spaces-groupoids-]{spaces-groupoids}
\externaldocument[spaces-more-groupoids-]{spaces-more-groupoids}
\externaldocument[bootstrap-]{bootstrap}
\externaldocument[spaces-pushouts-]{spaces-pushouts}
\externaldocument[spaces-chow-]{spaces-chow}
\externaldocument[groupoids-quotients-]{groupoids-quotients}
\externaldocument[spaces-more-cohomology-]{spaces-more-cohomology}
\externaldocument[spaces-simplicial-]{spaces-simplicial}
\externaldocument[spaces-duality-]{spaces-duality}
\externaldocument[formal-spaces-]{formal-spaces}
\externaldocument[restricted-]{restricted}
\externaldocument[spaces-resolve-]{spaces-resolve}
\externaldocument[formal-defos-]{formal-defos}
\externaldocument[defos-]{defos}
\externaldocument[cotangent-]{cotangent}
\externaldocument[examples-defos-]{examples-defos}
\externaldocument[algebraic-]{algebraic}
\externaldocument[examples-stacks-]{examples-stacks}
\externaldocument[stacks-sheaves-]{stacks-sheaves}
\externaldocument[criteria-]{criteria}
\externaldocument[artin-]{artin}
\externaldocument[quot-]{quot}
\externaldocument[stacks-properties-]{stacks-properties}
\externaldocument[stacks-morphisms-]{stacks-morphisms}
\externaldocument[stacks-limits-]{stacks-limits}
\externaldocument[stacks-cohomology-]{stacks-cohomology}
\externaldocument[stacks-perfect-]{stacks-perfect}
\externaldocument[stacks-introduction-]{stacks-introduction}
\externaldocument[stacks-more-morphisms-]{stacks-more-morphisms}
\externaldocument[stacks-geometry-]{stacks-geometry}
\externaldocument[moduli-]{moduli}
\externaldocument[moduli-curves-]{moduli-curves}
\externaldocument[examples-]{examples}
\externaldocument[exercises-]{exercises}
\externaldocument[guide-]{guide}
\externaldocument[desirables-]{desirables}
\externaldocument[coding-]{coding}
\externaldocument[obsolete-]{obsolete}
\externaldocument[fdl-]{fdl}
\externaldocument[index-]{index}

% Theorem environments.
%
\theoremstyle{plain}
\newtheorem{theorem}[subsection]{Theorem}
\newtheorem{proposition}[subsection]{Proposition}
\newtheorem{lemma}[subsection]{Lemma}

\theoremstyle{definition}
\newtheorem{definition}[subsection]{Definition}
\newtheorem{example}[subsection]{Example}
\newtheorem{exercise}[subsection]{Exercise}
\newtheorem{situation}[subsection]{Situation}

\theoremstyle{remark}
\newtheorem{remark}[subsection]{Remark}
\newtheorem{remarks}[subsection]{Remarks}

\numberwithin{equation}{subsection}

% Macros
%
\def\lim{\mathop{\mathrm{lim}}\nolimits}
\def\colim{\mathop{\mathrm{colim}}\nolimits}
\def\Spec{\mathop{\mathrm{Spec}}}
\def\Hom{\mathop{\mathrm{Hom}}\nolimits}
\def\Ext{\mathop{\mathrm{Ext}}\nolimits}
\def\SheafHom{\mathop{\mathcal{H}\!\mathit{om}}\nolimits}
\def\SheafExt{\mathop{\mathcal{E}\!\mathit{xt}}\nolimits}
\def\Sch{\mathit{Sch}}
\def\Mor{\mathop{\mathrm{Mor}}\nolimits}
\def\Ob{\mathop{\mathrm{Ob}}\nolimits}
\def\Sh{\mathop{\mathit{Sh}}\nolimits}
\def\NL{\mathop{N\!L}\nolimits}
\def\CH{\mathop{\mathrm{CH}}\nolimits}
\def\proetale{{pro\text{-}\acute{e}tale}}
\def\etale{{\acute{e}tale}}
\def\QCoh{\mathit{QCoh}}
\def\Ker{\mathop{\mathrm{Ker}}}
\def\Im{\mathop{\mathrm{Im}}}
\def\Coker{\mathop{\mathrm{Coker}}}
\def\Coim{\mathop{\mathrm{Coim}}}

% Boxtimes
%
\DeclareMathSymbol{\boxtimes}{\mathbin}{AMSa}{"02}

%
% Macros for moduli stacks/spaces
%
\def\QCohstack{\mathcal{QC}\!\mathit{oh}}
\def\Cohstack{\mathcal{C}\!\mathit{oh}}
\def\Spacesstack{\mathcal{S}\!\mathit{paces}}
\def\Quotfunctor{\mathrm{Quot}}
\def\Hilbfunctor{\mathrm{Hilb}}
\def\Curvesstack{\mathcal{C}\!\mathit{urves}}
\def\Polarizedstack{\mathcal{P}\!\mathit{olarized}}
\def\Complexesstack{\mathcal{C}\!\mathit{omplexes}}
% \Pic is the operator that assigns to X its picard group, usage \Pic(X)
% \Picardstack_{X/B} denotes the Picard stack of X over B
% \Picardfunctor_{X/B} denotes the Picard functor of X over B
\def\Pic{\mathop{\mathrm{Pic}}\nolimits}
\def\Picardstack{\mathcal{P}\!\mathit{ic}}
\def\Picardfunctor{\mathrm{Pic}}
\def\Deformationcategory{\mathcal{D}\!\mathit{ef}}
\usepackage{cleveref}
\crefname{originallemma}{original lemma}{original lemmas}
\Crefname{originallemma}{Original lemma}{Original lemmas}
\makeatletter
\newlabel{categories-lemma-functorial-limit}{{002L}{source}{}{sourcepointer-1}{}}
\newlabel{categories-lemma-functorial-limit@cref}{{[originallemma][0][]002L}{[1][0][]source}{}{}{}}
\makeatother
\crefname{originallemma}{original lemma}{original lemmas}
\Crefname{originallemma}{Original lemma}{Original lemmas}
\makeatletter
\newlabel{categories-lemma-functorial-colimit}{{002K}{source}{}{sourcepointer-2}{}}
\newlabel{categories-lemma-functorial-colimit@cref}{{[originallemma][0][]002K}{[1][0][]source}{}{}{}}
\makeatother
\title{The normalised complex of an Eilenberg-MacLane object is concentrated in one degree}
\author{The Stacks Project Authors}
\date{stacks-project:simplicial.tex:local-snapshot-20261012; GNU FDL 1.2 or later}
\begin{document}
\maketitle
\noindent\textit{Editorial note.} The counted statement is the author's own lemma environment labelled lemma-N-K (source lines 4012-4019); the proof environment begins on the second line after it and is its own proof, closing at source line 4030.\par\medskip
\noindent\textit{Editorial note.} The proof invokes two author results of the same chapter by label; both are retained as context spans and the generator expands the context transitively (budget 800 lines for this card), while reference targets in other chapters that carry canonical project tags are auto-declared as non-load-bearing source pointers.\par\medskip
\noindent\textit{Editorial note.} The retained preamble is the chapter's own minimal preamble (sources/stacks-simplicial/preamble.tex, taken with its real line count), a verbatim excerpt of the native Xy-pic setup, the native theorem declarations with the mathematical macros, and the project's own no-op reference environments; no macro was invented or redefined.\par\medskip
\noindent\textit{Editorial scope.} One original result of the Stacks Project's Simplicial Methods chapter, counted once: the normalised complex of K(A,k) is A concentrated in degree k.\par\medskip
\noindent\textit{Allowed standard prerequisites (editorial).}\begin{itemize}
\item Simplicial objects and the normalised complex. For a simplicial object of an abelian category, the normalised chain complex N, whose n-th term is the quotient of the n-simplices by the degenerate ones.. It is the object computed in the statement (source lines 4014-4018).
\item Eilenberg-MacLane objects. The simplicial object K(A,k) built as the induced object from the skeleton, with K(A,k)\_i = 0 for i below k, equal to A in degree k, and nondegenerate simplices only in degree k (the author's definition-eilenberg-maclane and lemma-imshriek-abelian, retained as context).. It supplies the three degree computations of the proof (source lines 4022-4029).
\end{itemize}\medskip
\ifdefined\hypertarget\hypertarget{sourcepointer-1}{}\fi
\noindent\textit{Source reference.} Original lemma 002L: sources/stacks-chow/tags.txt, source lines 102--102. Cross-chapter pointer whose target carries the canonical tag 002L in the project tag list; declared so the reference resolves inside the assembled package, while every step the counted argument uses remains inside the retained spans.\par\medskip
\ifdefined\hypertarget\hypertarget{sourcepointer-2}{}\fi
\noindent\textit{Source reference.} Original lemma 002K: sources/stacks-chow/tags.txt, source lines 101--101. Cross-chapter pointer whose target carries the canonical tag 002K in the project tag list; declared so the reference resolves inside the assembled package, while every step the counted argument uses remains inside the retained spans.\par\medskip
\begin{lemma}
\label{lemma-imshriek-abelian}
Let $\mathcal{A}$ be an abelian category
Let $U$ be an $m$-truncated simplicial object of
$\mathcal{A}$. For $n > m$ we have $N(i_{m!}U)_n = 0$.
\end{lemma}

\begin{proof}
Write $V = i_{m!}U$. Let $V' \subset V$ be the
simplicial subobject of $V$ with $V'_i = V_i$ for $i \leq m$
and $N(V'_i) = 0$ for $i > m$,
see Lemma \ref{lemma-simplicial-abelian-n-skel-sub}.
By the adjunction formula,
since $\text{sk}_m V' = U$, there is an inverse to the
injection $V' \to V$. Hence $V' = V$.
\end{proof}

\begin{definition}
\label{definition-eilenberg-maclane}
Let $\mathcal{A}$ be an abelian category.
Let $A$ be an object of $\mathcal{A}$ and
let $k$ be an integer $\geq 0$.
The {\it Eilenberg-Maclane object $K(A, k)$}
is given by the object $K(A, k) = i_{k!}U$
which is described in
Lemma \ref{lemma-eilenberg-maclane-object} above.
\end{definition}

\begin{lemma}
\label{lemma-simplicial-abelian-n-skel-sub}
Let $\mathcal{A}$ be an abelian category.
Let $U$ be a simplicial object of $\mathcal{A}$.
Let $n \geq 0$ be an integer.
The rule
$$
U'_m = \sum\nolimits_{\varphi : [m] \to [i], \ i\leq n} \Im(U(\varphi))
$$
defines a sub simplicial object $U' \subset U$ with $U'_i = U_i$
for $i \leq n$.
Moreover, $N(U'_m) = 0$ for all $m > n$.
\end{lemma}

\begin{proof}
Pick $m$, $i \leq n$ and some $\varphi : [m] \to [i]$.
The image under $U(\psi)$ of $\Im(U(\varphi))$
for any $\psi : [m'] \to [m]$ is
equal to the image of $U(\varphi \circ \psi)$ and
$\varphi \circ \psi : [m'] \to [i]$.
Hence $U'$ is a simplicial object.
Pick $m > n$. We have to show $N(U'_m) = 0$.
By definition of $N(U_m)$ and $N(U'_m)$ we have
$N(U'_m) = U'_m \cap N(U_m)$ (intersection of subobjects).
Since $U$ is split by Lemma \ref{lemma-splitting-abelian-category},
it suffices to show that $U'_m$ is contained in the sum
$$
\sum\nolimits_{\varphi : [m] \to [m']\text{ surjective}, \ m' < m}
\Im(U(\varphi)|_{N(U_{m'})}).
$$
By the splitting each $U_{m'}$ is the sum of images of
$N(U_{m''})$ via $U(\psi)$ for surjective maps
$\psi : [m'] \to [m'']$. Hence the displayed sum above
is the same as
$$
\sum\nolimits_{\varphi : [m] \to [m']\text{ surjective}, \ m' < m}
\Im(U(\varphi)).
$$
Clearly $U'_m$ is contained in this by the simple fact that
any $\varphi : [m] \to [i]$, $i \leq n$ occurring in the definition
of $U'_m$ may be factored as
$[m] \to [m'] \to [i]$ with $[m] \to [m']$ surjective
and $m' < m$ as in the last displayed sum above.
\end{proof}

\begin{lemma}
\label{lemma-eilenberg-maclane-object}
With $A$, $k$ and $U$ as above, so $U_i = 0$, $i < k$ and $U_k = A$.
\begin{enumerate}
\item Given a $k$-truncated simplicial object $V$
we have
$$
\Mor(U, V)
=
\{ f : A \to V_k \mid d^k_i \circ f = 0, \ i = 0, \ldots, k \}
$$
and
$$
\Mor(V, U)
=
\{ f : V_k \to A \mid f \circ s^{k - 1}_i = 0, \ i = 0, \ldots, k - 1 \}.
$$
\item The object $i_{k!} U$ has $n$th term equal to
$\bigoplus_\alpha A$ where $\alpha$ runs over all
surjective morphisms $\alpha : [n] \to [k]$.
\item For any $\varphi : [m] \to [n]$ the map
$i_{k!} U(\varphi)$ is described as the mapping
$\bigoplus_\alpha A \to \bigoplus_{\alpha'} A$
which maps to component corresponding to $\alpha : [n] \to [k]$
to zero if $\alpha \circ \varphi$ is not surjective and
by the identity to the component corresponding to
$\alpha \circ \varphi$ if it is surjective.
\item The object $\text{cosk}_k U$ has $n$th term equal to
$\bigoplus_\beta A$, where $\beta$ runs over all
injective morphisms $\beta : [k] \to [n]$.
\item For any $\varphi : [m] \to [n]$ the map
$\text{cosk}_k U(\varphi)$ is described as the mapping
$\bigoplus_\beta A \to \bigoplus_{\beta'} A$
which maps to component corresponding to $\beta : [k] \to [n]$
to zero if $\beta$ does not factor through $\varphi$ and
by the identity to each of the components corresponding to
$\beta'$ such that $\beta = \varphi \circ \beta'$
if it does.
\item The canonical map
$
c : i_{k !} U \to \text{cosk}_k U
$
in degree $n$ has $(\alpha, \beta)$ coefficient $A \to A$
equal to zero if $\alpha \circ \beta$ is not the identity
and equal to $\text{id}_A$ if it is.
\item The canonical map
$
c : i_{k !} U \to \text{cosk}_k U
$
is injective.
\end{enumerate}
\end{lemma}

\begin{proof}
The proof of (1) is left to the reader.

\medskip\noindent
Let us take the rules of (2) and (3)
as the definition of a simplicial object, call it $\tilde U$.
We will show that it is an incarnation of $i_{k!}U$.
This will prove (2), (3) at the same time. We have to show
that given a morphism $f : U \to \text{sk}_kV$
there exists a unique morphism $\tilde f : \tilde U \to V$
which recovers $f$ upon taking the $k$-skeleton.
From (1) we see that $f$ corresponds with a morphism
$f_k : A \to V_k$ which maps into the kernel of
$d^k_i$ for all $i$. For any surjective $\alpha : [n] \to [k]$
we set $\tilde f_\alpha : A \to V_n$ equal to the composition
$\tilde f_\alpha = V(\alpha) \circ f_k : A \to V_n$. We define
$\tilde f_n : \tilde U_n \to V_n$ as the sum of
the $\tilde f_\alpha$ over $\alpha : [n] \to [k]$ surjective.
Such a collection of $\tilde f_\alpha$ defines a morphism
of simplicial objects if and only if
for any $\varphi : [m] \to [n]$ the diagram
$$
\xymatrix{
\bigoplus_{\alpha : [n] \to [k]\text{ surjective}} A
\ar[r]_-{\tilde f_n}
\ar[d]_{(3)} &
V_n \ar[d]^{V(\varphi)} \\
\bigoplus_{\alpha' : [m] \to [k]\text{ surjective}} A
\ar[r]^-{\tilde f_m} &
V_m
}
$$
is commutative. Choosing $\varphi = \alpha$ shows our choice of
$\tilde f_\alpha$ is uniquely determined by $f_k$.
The commutativity in general may be checked for each summand
of the left upper corner separately. It is clear for the
summands corresponding to $\alpha$ where
$\alpha \circ \varphi$ is surjective, because those get
mapped by $\text{id}_A$ to the summand with
$\alpha' = \alpha \circ \varphi$, and we have
$\tilde f_{\alpha'} = V(\alpha') \circ f_k =
V(\alpha \circ \varphi) \circ f_k = V(\varphi) \circ \tilde f_\alpha$.
For those where $\alpha \circ \varphi$
is not surjective, we have to show that $V(\varphi) \circ \tilde f_\alpha = 0$.
By definition this is equal to
$V(\varphi) \circ V(\alpha) \circ f_k = V(\alpha \circ \varphi) \circ f_k$.
Since $\alpha \circ \varphi$ is not surjective we can write it
as $\delta^k_i \circ \psi$, and we deduce that
$V(\varphi) \circ V(\alpha) \circ f_k =
V(\psi) \circ d^k_i \circ f_k = 0$ see above.

\medskip\noindent
Let us take the rules of (4) and (5)
as the definition of a simplicial object, call it $\tilde U$.
We will show that it is an incarnation of $\text{cosk}_k U$.
This will prove (4), (5) at the same time. The argument is completely dual
to the proof of (2), (3) above, but we give it anyway.
We have to show
that given a morphism $f : \text{sk}_kV \to U$
there exists a unique morphism $\tilde f : V \to \tilde U$
which recovers $f$ upon taking the $k$-skeleton.
From (1) we see that $f$ corresponds with a morphism
$f_k : V_k \to A$ which is zero on the image of $s^{k - 1}_i$
for all $i$. For any injective $\beta : [k] \to [n]$
we set $\tilde f_\beta : V_n \to A$ equal to the composition
$\tilde f_\beta = f_k \circ V(\beta) : V_n \to A$. We define
$\tilde f_n : V_n \to \tilde U_n$ as the sum of
the $\tilde f_\beta$ over $\beta : [k] \to [n]$ injective.
Such a collection of $\tilde f_\beta$ defines a morphism
of simplicial objects if and only if
for any $\varphi : [m] \to [n]$ the diagram
$$
\xymatrix{
V_n
\ar[d]_{V(\varphi)}
\ar[r]_-{\tilde f_n}
&
\bigoplus_{\beta : [k] \to [n]\text{ injective}} A
\ar[d]^{(5)}
\\
V_m
\ar[r]^-{\tilde f_m}
&
\bigoplus_{\beta' : [k] \to [m]\text{ injective}} A
}
$$
is commutative. Choosing $\varphi = \beta$ shows our choice of
$\tilde f_\beta$ is uniquely determined by $f_k$.
The commutativity in general may be checked for each summand
of the right lower corner separately. It is clear for the
summands corresponding to $\beta'$ where
$\varphi \circ \beta'$ is injective, because these summands
get mapped into by exactly the summand with
$\beta = \varphi \circ \beta'$ and we have in that case
$\tilde f_{\beta'} \circ V(\varphi) =
f_k \circ V(\beta') \circ V(\varphi) =
f_k \circ V(\beta) = \tilde f_\beta$. For those where
$\varphi \circ \beta'$ is not injective,
we have to show that $\tilde f_{\beta'} \circ V(\varphi) = 0$.
By definition this is equal to
$f_k \circ V(\beta') \circ V(\varphi) =
f_k \circ V(\varphi \circ \beta')$.
Since $\varphi \circ \beta'$ is not injective we can write it
as $\psi \circ \sigma^{k - 1}_i$, and we deduce that
$f_k \circ V(\beta') \circ V(\varphi) =
f_k \circ s^{k - 1}_i \circ V(\psi) = 0$ see above.

\medskip\noindent
The composition $i_{k!}U \to \text{cosk}_kU$ is the
unique map of simplicial objects which is
the identity on $A = U_k = (i_{k!}U)_k = (\text{cosk}_kU)_k$.
Hence it suffices to check that the proposed rule defines
a morphism of simplicial objects.
To see this we have to show that
for any $\varphi : [m] \to [n]$ the diagram
$$
\xymatrix{
\bigoplus_{\alpha : [n] \to [k]\text{ surjective}} A
\ar[d]_{(3)}
\ar[r]_{(6)}
&
\bigoplus_{\beta : [k] \to [n]\text{ injective}} A
\ar[d]^{(5)}
\\
\bigoplus_{\alpha' : [m] \to [k]\text{ surjective}} A
\ar[r]^{(6)}
&
\bigoplus_{\beta' : [k] \to [m]\text{ injective}} A
}
$$
is commutative. Now we can think of this in terms of
matrices filled with only $0$'s and $1$'s as follows:
The matrix of (3) has a nonzero
$(\alpha', \alpha)$ entry if and only if
$\alpha' = \alpha \circ \varphi$. Likewise
the matrix of (5) has a nonzero
$(\beta', \beta)$ entry if and only if
$\beta = \varphi \circ \beta'$. The upper matrix
of (6) has a nonzero $(\alpha, \beta)$ entry if and only if
$\alpha \circ \beta = \text{id}_{[k]}$. Similarly for the
lower matrix of (6). The commutativity of the
diagram then comes down to computing the
$(\alpha, \beta')$ entry for both compositions
and seeing they are equal. This comes down to the
following equality
$$
\# \left\{
\beta
\mid
\beta = \varphi \circ \beta'
\text{ and }
\alpha \circ \beta = \text{id}_{[k]}
\right\}
=
\# \left\{
\alpha'
\mid
\alpha' = \alpha \circ \varphi
\text{ and }
\alpha' \circ \beta' = \text{id}_{[k]}
\right\}
$$
whose proof may safely be left to the reader.

\medskip\noindent
Finally, we prove (7). This follows directly from
Lemmas \ref{lemma-injective-map-simplicial-abelian},
\ref{lemma-recover-cosk}, \ref{lemma-recovering-U-for-real}
and \ref{lemma-imshriek-abelian}.
\end{proof}

\begin{lemma}
\label{lemma-splitting-abelian-category}
Let $\mathcal{A}$ be an abelian category.
Let $U$ be a simplicial object in $\mathcal{A}$.
Then $U$ has a splitting obtained by taking $N(U_0) = U_0$ and
for $m \geq 1$ taking
$$
N(U_m) = \bigcap\nolimits_{i = 0}^{m - 1} \Ker(d^m_i).
$$
Moreover, this splitting is functorial on the category of
simplicial objects of $\mathcal{A}$.
\end{lemma}

\begin{proof}
For any object $A$ of $\mathcal{A}$ we obtain
a simplicial abelian group $\Mor_\mathcal{A}(A, U)$.
Each of these are canonically split by Lemma
\ref{lemma-splitting-simplicial-groups}. Moreover,
$$
N(\Mor_\mathcal{A}(A, U_m)) =
\bigcap\nolimits_{i = 0}^{m - 1} \Ker(d^m_i) =
\Mor_\mathcal{A}(A, N(U_m)).
$$
Hence we see that the morphism (\ref{equation-splitting})
becomes an isomorphism after applying the functor
$\Mor_\mathcal{A}(A, -)$ for any object of $\mathcal{A}$.
Hence it is an isomorphism by the Yoneda lemma.
\end{proof}

\begin{lemma}
\label{lemma-injective-map-simplicial-abelian}
\begin{slogan}
The Dold-Kan normalization functor reflects
injectivity, surjectivity, and isomorphy.
\end{slogan}
Let $\mathcal{A}$ be an abelian category.
Let $f : U \to V$ be a morphism of
simplicial objects of $\mathcal{A}$.
If the induced morphisms $N(f)_i : N(U)_i \to N(V)_i$
are injective for all $i$, then $f_i$ is
injective for all $i$. Same holds with ``injective'' replaced
with ``surjective'', or ``isomorphism''.
\end{lemma}

\begin{proof}
This is clear from Lemma \ref{lemma-splitting-abelian-category}
and the definition of a splitting.
\end{proof}

\begin{lemma}
\label{lemma-recover-cosk}
Let $\mathcal{C}$ be a category with finite limits.
Let $U$ be an $n$-truncated simplicial object of $\mathcal{C}$.
The morphism $\text{sk}_n \text{cosk}_n U \to U$
is an isomorphism.
\end{lemma}

\begin{proof}
Combine Lemmas \ref{lemma-existence-cosk} and \ref{lemma-trivial-cosk}.
\end{proof}

\begin{lemma}
\label{lemma-recovering-U-for-real}
Let $\mathcal{C}$ be a category which has finite colimits.
Let $U$ be an $m$-truncated simplicial object of $\mathcal{C}$.
The map $U \to \text{sk}_m i_{m!}U$ is an isomorphism.
\end{lemma}

\begin{proof}
Combine Lemmas \ref{lemma-left-adjoint-exists} and \ref{lemma-recovering-U}.
\end{proof}

\begin{equation}
\label{equation-splitting}
\coprod\nolimits_{\varphi : [n] \to [m]\text{ surjective}}
N(U_m)
\longrightarrow
U_n
\end{equation}

\begin{lemma}
\label{lemma-splitting-simplicial-groups}
Let $U$ be a simplicial abelian group.
Then $U$ has a splitting obtained by taking $N(U_0) = U_0$ and
for $m \geq 1$ taking
$$
N(U_m) = \bigcap\nolimits_{i = 0}^{m - 1} \Ker(d^m_i).
$$
Moreover, this splitting is functorial on the category
of simplicial abelian groups.
\end{lemma}

\begin{proof}
By induction on $n$ we will show that the choice of $N(U_m)$
in the lemma guarantees that (\ref{equation-splitting}) is
an isomorphism for $m \leq n$. This is clear for $n = 0$.
In the rest of this proof we are going to
drop the superscripts from the maps $d_i$ and $s_i$ in order
to improve readability. We will also repeatedly use the relations
from Remark \ref{remark-relations}.

\medskip\noindent
First we make a general remark.
For $0 \leq i \leq m$ and $z \in U_m$ we have
$d_i(s_i(z)) = z$. Hence we can write
any $x \in U_{m + 1}$ uniquely as
$x = x' + x''$ with $d_i(x') = 0$
and $x'' \in \Im(s_i)$
by taking $x' = (x - s_i(d_i(x)))$ and
$x'' = s_i(d_i(x))$. Moreover, the element
$z \in U_m$ such that $x'' = s_i(z)$
is unique because $s_i$ is injective.

\medskip\noindent
Here is a procedure for decomposing any $x \in U_{n + 1}$.
First, write $x = x_0 + s_0(z_0)$ with $d_0(x_0) = 0$.
Next, write $x_0 = x_1 + s_1(z_1)$ with
$d_1(x_1) = 0$. Continue like this to get
\begin{eqnarray*}
x & = & x_0 + s_0(z_0), \\
x_0 & = & x_1 + s_1(z_1), \\
x_1 & = & x_2 + s_2(z_2), \\
\ldots & \ldots & \ldots \\
x_{n - 1} & = & x_n + s_n(z_n)
\end{eqnarray*}
where $d_i(x_i) = 0$ for all $i = n, \ldots, 0$.
By our general remark above all of the $x_i$
and $z_i$ are determined uniquely by $x$.
We claim that
$x_i \in
\Ker(d_0) \cap
\Ker(d_1) \cap
\ldots \cap
\Ker(d_i)$
and
$z_i \in
\Ker(d_0) \cap
\ldots \cap
\Ker(d_{i - 1})$
for $i = n, \ldots, 0$.
Here and in the following
an empty intersection of kernels indicates
the whole space; i.e.,
the notation
$z_0 \in \Ker(d_0) \cap
\ldots \cap
\Ker(d_{i - 1})$
when $i = 0$ means $z_0 \in U_n$ with no restriction.

\medskip\noindent
We prove this by ascending induction on $i$.
It is clear for $i = 0$ by construction of $x_0$ and $z_0$.
Let us prove it for $0 < i \leq n$ assuming the result for $i - 1$.
First of all we have $d_i(x_i) = 0$ by construction.
So pick a $j$ with $0 \leq j < i$. We have
$d_j(x_{i - 1}) = 0$ by induction. Hence
$$
0 = d_j(x_{i - 1})
= d_j(x_i) + d_j(s_i(z_i))
= d_j(x_i) + s_{i - 1}(d_j(z_i)).
$$
The last equality by the relations of Remark \ref{remark-relations}.
These relations also imply that
$d_{i - 1}(d_j(x_i)) = d_j(d_i(x_i)) = 0$
because $d_i(x_i)= 0$ by construction.
Then the uniqueness in the general remark above shows the equality
$0 = x' + x'' = d_j(x_i) + s_{i - 1}(d_j(z_i))$
can only hold if both terms are zero. We conclude that
$d_j(x_i) = 0$ and by injectivity of $s_{i - 1}$ we also
conclude that $d_j(z_i) = 0$. This proves the claim.

\medskip\noindent
The claim implies we can uniquely write
$$
x = s_0(z_0) + s_1(z_1) + \ldots + s_n(z_n) + x_n
$$
with $x_n \in N(U_{n + 1})$ and
$z_i \in \Ker(d_0) \cap \ldots \cap \Ker(d_{i - 1})$.
We can reformulate this as saying that we have found a direct
sum decomposition
$$
U_{n + 1}
=
N(U_{n + 1})
\oplus
\bigoplus\nolimits_{i = 0}^{i = n}
s_i\Big(\Ker(d_0) \cap \ldots \cap \Ker(d_{i - 1})\Big)
$$
with the property that
$$
\Ker(d_0) \cap \ldots \cap \Ker(d_j)
=
N(U_{n + 1}) \oplus
\bigoplus\nolimits_{i = j + 1}^{i = n}
s_i\Big(\Ker(d_0) \cap \ldots \cap \Ker(d_{i - 1})\Big)
$$
for $j = 0, \ldots, n$.
The result follows from this statement as follows.
Each of the $z_i$ in the expression for $x$
can be written uniquely as
$$
z_i = s_i(z'_{i, i}) + \ldots + s_{n - 1}(z'_{i, n - 1}) + z_{i, 0}
$$
with $z_{i, 0} \in N(U_n)$ and
$z'_{i, j} \in \Ker(d_0) \cap \ldots \cap \Ker(d_{j - 1})$.
The first few steps in the decomposition of $z_i$ are zero because
$z_i$ already is in the kernel of $d_0, \ldots, d_i$.
This in turn uniquely gives
$$
x = x_n + s_0(z_{0, 0}) + s_1(z_{1, 0}) + \ldots + s_n(z_{n, 0}) +
\sum\nolimits_{0 \leq i \leq j \leq n - 1} s_i(s_j(z'_{i, j})).
$$
Continuing in this fashion we see that we in the end obtain
a decomposition of $x$ as a sum of terms
of the form
$$
s_{i_1} s_{i_2} \ldots s_{i_k} (z)
$$
with $0 \leq i_1 \leq i_2 \leq \ldots \leq i_k \leq n - k + 1$ and
$z \in N(U_{n + 1 - k})$. This is exactly the required
decomposition, because any surjective map $[n + 1] \to [n + 1 - k]$
can be uniquely expressed in the form
$$
\sigma^{n - k}_{i_k} \ldots \sigma^{n - 1}_{i_2} \sigma^n_{i_1}
$$
with $0 \leq i_1 \leq i_2 \leq \ldots \leq i_k \leq n - k + 1$.
\end{proof}

\begin{lemma}
\label{lemma-existence-cosk}
If the category $\mathcal{C}$ has finite limits, then
$\text{cosk}_m$ functors exist for all $m$. Moreover,
for any $m$-truncated simplicial object $U$ the
simplicial object $\text{cosk}_mU$ is described
by the formula
$$
(\text{cosk}_mU)_n = \lim_{(\Delta/[n])_{\leq m}^{opp}} U(n)
$$
and for $\varphi : [n] \to [n']$ the map
$\text{cosk}_mU(\varphi)$ comes from the
identification $U(n') \circ \overline{\varphi} = U(n)$ above
via Categories, Lemma \ref{categories-lemma-functorial-limit}.
\end{lemma}

\begin{proof}
During the proof of this lemma we denote $\text{cosk}_mU$ the
simplicial object with $(\text{cosk}_mU)_n$ equal to
$\lim_{(\Delta/[n])_{\leq m}^{opp}} U(n)$.
We will conclude at the end of the proof that it does
satisfy the required mapping property.

\medskip\noindent
Suppose that $V$ is a simplicial object.
A morphism $\gamma : V \to \text{cosk}_mU$ is given by a sequence
of morphisms $\gamma_n : V_n \to (\text{cosk}_mU)_n$.
By definition of a limit, this is given by a
collection of morphisms $\gamma(\alpha) : V_n \to U_k$
where $\alpha$ ranges over all $\alpha : [k] \to [n]$
with $k \leq m$. These morphisms then also satisfy
the rules that
$$
\xymatrix{
V_n \ar[r]_{\gamma(\alpha)} &  U_k \\
V_{n'} \ar[r]^{\gamma(\alpha')} \ar[u]^{V(\varphi)} & U_{k'} \ar[u]_{U(\psi)}
}
$$
are commutative, given any $0 \leq k, k' \leq m$, $0 \leq n, n'$
and any $\psi : [k] \to [k']$, $\varphi : [n] \to [n']$,
$\alpha : [k] \to [n]$ and $\alpha' : [k'] \to [n']$ in $\Delta$
such that $\varphi \circ \alpha = \alpha' \circ \psi$.
Taking $n = k = k'$, $\varphi = \alpha'$, and $\alpha = \psi = \text{id}_{[k]}$
we deduce that $\gamma(\alpha') = \gamma(\text{id}_{[k]}) \circ V(\alpha')$.
In other words, the morphisms $\gamma(\text{id}_{[k]})$, $k \leq m$
determine the morphism $\gamma$. And it is easy to see that these
morphisms form a morphism $\text{sk}_m V \to U$.

\medskip\noindent
Conversely, given a morphism $\gamma : \text{sk}_m V \to U$,
we obtain a family of morphisms $\gamma(\alpha)$
where $\alpha$ ranges over all $\alpha : [k] \to [n]$
with $k \leq m$ by setting $\gamma(\alpha) =
\gamma(\text{id}_{[k]}) \circ V(\alpha)$. These morphisms
satisfy all the displayed commutativity restraints pictured
above, and hence give rise to a morphism $V \to \text{cosk}_m U$.
\end{proof}

\begin{lemma}
\label{lemma-trivial-cosk}
Let $\mathcal{C}$ be a category.
Let $U$ be an $m$-truncated simplicial object of $\mathcal{C}$.
For $n \leq m$ the limit $\lim_{(\Delta/[n])_{\leq m}^{opp}} U(n)$
exists and is canonically isomorphic to $U_n$.
\end{lemma}

\begin{proof}
This is true because the category $(\Delta/[n])_{\leq m}$
has an final object in this case, namely the identity
map $[n] \to [n]$.
\end{proof}

\begin{lemma}
\label{lemma-left-adjoint-exists}
Let $\mathcal{C}$ be a category which has finite colimits.
The functors $i_{m!}$ exist for all $m$.
Let $U$ be an $m$-truncated simplicial object of $\mathcal{C}$.
The simplicial object $i_{m!}U$
is described by the formula
$$
(i_{m!}U)_n = \colim_{([n]/\Delta)_{\leq m}^{opp}} U(n)
$$
and for $\varphi : [n] \to [n']$ the map
$i_{m!}U(\varphi)$ comes from the
identification $U(n) \circ \underline{\varphi} = U(n')$ above
via Categories, Lemma \ref{categories-lemma-functorial-colimit}.
\end{lemma}

\begin{proof}
In this proof we denote $i_{m!}U$ the simplicial object
whose $n$th term is given by the displayed formula of the
lemma. We will show it satisfies the adjointness property.

\medskip\noindent
Let $V$ be a simplicial object of $\mathcal{C}$.
Let $\gamma : U \to \text{sk}_mV$ be given.
A morphism
$$
\colim_{([n]/\Delta)_{\leq m}^{opp}} U(n) \to T
$$
is given by a compatible system of morphisms
$f_\alpha : U_k \to T$ where $\alpha : [n] \to [k]$
with $k \leq m$. Certainly, we have such a system of
morphisms by taking the compositions
$$
U_k \xrightarrow{\gamma_k} V_k \xrightarrow{V(\alpha)} V_n.
$$
Hence we get an induced morphism $(i_{m!}U)_n \to V_n$.
We leave it to the reader to see that these form a
morphism of simplicial objects $\gamma' : i_{m!}U \to V$.

\medskip\noindent
Conversely, given a morphism $\gamma' : i_{m!}U \to V$ we obtain
a morphism $\gamma : U \to \text{sk}_m V$ by setting
$\gamma_i : U_i \to V_i$ equal to the composition
$$
U_i
\xrightarrow{\text{id}_{[i]}}
\colim_{([i]/\Delta)_{\leq m}^{opp}} U(i)
\xrightarrow{\gamma'_i}
V_i
$$
for $0 \leq i \leq n$. We leave it to the reader to see that
this is the inverse of the construction above.
\end{proof}

\begin{lemma}
\label{lemma-recovering-U}
Let $\mathcal{C}$ be a category.
Let $U$ be an $m$-truncated simplicial object of $\mathcal{C}$.
For any $n \leq m$ the colimit
$$
\colim_{([n]/\Delta)_{\leq m}^{opp}} U(n)
$$
exists and is equal to $U_n$.
\end{lemma}

\begin{proof}
This is so because the category $([n]/\Delta)_{\leq m}$
has an initial object, namely $\text{id} : [n] \to [n]$.
\end{proof}

\begin{remark}
\label{remark-relations}
By abuse of notation we sometimes write $d_i : U_n \to U_{n - 1}$
instead of $d^n_i$, and similarly for $s_i : U_n \to U_{n + 1}$.
The relations among the morphisms $d^n_i$ and $s^n_i$
may be expressed as follows:
\begin{enumerate}
\item If $i < j$, then $d_i \circ d_j = d_{j - 1} \circ d_i$.
\item If $i < j$, then $d_i \circ s_j = s_{j - 1} \circ d_i$.
\item We have $\text{id} = d_j \circ s_j = d_{j + 1} \circ s_j$.
\item If $i > j + 1$, then $d_i \circ s_j = s_j \circ d_{i - 1}$.
\item If $i \leq j$, then $s_i \circ s_j = s_{j + 1} \circ s_i$.
\end{enumerate}
This means that whenever the compositions on both the left and the
right are defined then the corresponding equality should hold.
\end{remark}

\begin{lemma}
\label{lemma-N-K}
Let $\mathcal{A}$ be an abelian category.
Let $A$ be an object of $\mathcal{A}$ and
let $k$ be an integer. We have
$N(K(A, k))_i = A$ if $i = k$ and
$0$ else.
\end{lemma}

\begin{proof}
It is clear that $N(K(A, k))_i = 0$ when $i < k$
because $K(A, k)_i = 0$ in that case.
It is clear that $N(K(A, k))_k = A$ since
$K(A, k)_{k - 1} = 0$ and $K(A, k)_k = A$.
For $i > k$ we have $N(K(A, k))_i = 0$
by Lemma \ref{lemma-imshriek-abelian} and
the definition of $K(A, k)$, see Definition
\ref{definition-eilenberg-maclane}.
\end{proof}

\end{document}
